\documentclass[11pt,a4paper]{article}
\usepackage[margin=2.1cm]{geometry}
\usepackage[T1]{fontenc}
\usepackage{lmodern}
\usepackage{microtype}
\usepackage{booktabs}
\usepackage{tabularx}
\usepackage{array}
\usepackage{enumitem}
\usepackage{url}

\hyphenpenalty=10000
\exhyphenpenalty=10000
\emergencystretch=3em
\linespread{1.02}
\setlist{itemsep=3pt,topsep=4pt,parsep=0pt}
\setlength{\parindent}{0pt}
\setlength{\parskip}{4pt}
\renewcommand{\arraystretch}{1.15}

\title{RAG Assignment Report: Baseline Retrieval vs Cross Encoder Reranking}
\author{MD TAUFIQUE HUSSAIN \quad Roll No. 2025202007}
\date{GitHub: \url{https://github.com/Taufiquesir/rag-baseline-vs-reranking}}

\begin{document}
\maketitle
\vspace{-1.2em}

\section{Introduction and setup}

This report compares two ways of choosing the 3 chunks that are given to the LLM in a small RAG system built over one technical PDF. In Pipeline 1 (baseline) a Bi Encoder retrieves the top 10 chunks by cosine similarity and the first 3 go to the LLM, which writes Answer 1. In Pipeline 2 (reranking) the same 10 chunks are scored again by a Cross Encoder, sorted by that score, and the new top 3 go to the LLM, which writes Answer 2. The chunks, the top 10 candidates, the LLM and the prompt are identical in both pipelines, so any difference between the two answers comes only from reranking.

\begin{itemize}
\item \textbf{Document.} NIST SP 800-63B-4, Digital Identity Guidelines: Authentication and Authenticator Management (July 2025, 129 pages), the US standard for passwords, multifactor login, one time codes, passkeys and sessions. It was chosen because it should be hard for plain vector search: many rules use SHALL NOT, many rules have an exception (``unless'', ``however''), the paragraphs for assurance levels AAL1, AAL2 and AAL3 are nearly identical and differ mainly in a number, and the standard says ``PSTN'' where a user would say ``SMS''.
\item \textbf{Preprocessing.} The notebook removes the first 11 pages (title, legal notes, table of contents) and the header printed on every page, and then cuts the text into chunks of 120 words with 30 words of overlap. This gives 413 chunks.
\item \textbf{Models.} Bi Encoder \texttt{all-MiniLM-L6-v2} with exact FAISS cosine search, Cross Encoder \texttt{ms-marco-MiniLM-L-6-v2}, and the LLM Qwen2.5 1.5B Instruct on a laptop CPU with greedy decoding. The prompt tells the LLM to answer only from the given passages or reply ``Not found''.
\item \textbf{Checking retrieval and answers.} For each test question, the chunk that contains the exact sentence answering it is the gold chunk, and its rank was recorded for both pipelines. Only ranks 1 to 3 reach the LLM. Answers were judged by hand against the sentence in the standard, and an answer was counted as correct when it states that rule.
\end{itemize}

\section{Architecture comparison}

A Bi Encoder reads the question and each chunk separately. It turns each of them into one vector and compares the two vectors with cosine similarity. Because a chunk vector does not depend on the question, all 413 chunk vectors were computed once in advance (13.5 seconds), and encoding a question and searching all vectors takes only about 18 ms. The weakness is that a whole chunk is squeezed into one vector before the model has seen the question, so small details can be blurred: a number, a level name such as AAL1, or the word ``not''. This happened in query 6, which asks about AAL1. The AAL2 paragraph (chunk 26, cosine 0.752) and the AAL3 paragraph (chunk 31, 0.747) both scored higher than the AAL1 paragraph (chunk 20, 0.731), so the Bi Encoder did not separate the level that the question asked about.

A Cross Encoder reads the question and the chunk together as one input, and self attention lets every question word look at every chunk word. This is a likely reason why it could rank the AAL1 paragraph first, but the experiments do not isolate the cause. The price is that nothing can be prepared in advance: it has to run once for every question and chunk pair, which took about 37 ms per pair on the laptop used here. This is why the two models are combined. The fast Bi Encoder picks 10 candidates from the whole document and the slower Cross Encoder only orders those 10.

\section{Qualitative and rank shift analysis}

Reranking moved the gold chunk into the top 3 in all 8 failure queries, so retrieval improved in all eight cases. The final answers did not improve in all eight. With the reranked context the LLM gave a correct answer in 6 of the 8 queries (not in queries 2 and 4). With the baseline context none of the 8 answers was fully correct, and query 3 had the right conclusion with a wrong reason. The table gives the gold chunk rank and what the LLM said. Queries 9 to 11 are counter examples where reranking did not help.

{\small
\newcolumntype{R}[1]{>{\raggedright\arraybackslash}p{#1}}
\begin{tabularx}{\textwidth}{@{}r R{1.9cm} >{\raggedright\arraybackslash}X c c R{2.6cm} R{3.1cm}@{}}
\toprule
\# & Type & Question & Base & Rerank & Answer 1 & Answer 2 \\
\midrule
1 & Exception & When must a user be forced to change a password? & \#5 & \#1 & unclear, only points to a section & correct, only after compromise \\
2 & Exception & Is encryption required for one time codes sent by SMS? & \#10 & \#1 & not found & wrong, missed the PSTN exception \\
3 & Exception & Can synced passkeys be used at AAL3? & \#7 & \#3 & right answer, wrong reason & correct, keys must not be exportable \\
4 & Negation & Must users change passwords every 90 days? & \#5 & \#1 & not found & not found \\
5 & Negation & Does a new OTP reset the failed attempt counter? & \#5 & \#1 & not found & correct, SHALL NOT reset \\
6 & Keyword trap & What is the reauthentication timeout at AAL1? & \#6 & \#1 & not found & correct, 30 days \\
7 & Keyword trap & What is the inactivity timeout at AAL3? & \#4 & \#1 & wrong, 1 hour & correct, 15 minutes \\
8 & Synonym & Is the use of SMS for login restricted? & \#4 & \#1 & confused & correct, PSTN is restricted \\
9 & Counter & Can a fingerprint alone be used to log in? & \#3 & \#7 & no, wrong reason & no, wrong reason \\
10 & Counter & What is the inactivity timeout at AAL2? & \#2 & \#3 & correct & correct \\
11 & Counter & Can email be used to send one time login codes? & none & none & not found & hallucinated yes \\
\bottomrule
\end{tabularx}
}

\textbf{Clearest case, query 7.} The baseline top 3 were a general session section (chunk 218) and two chunks from the AAL2 text (chunks 25 and 26), which give the inactivity timeout as no more than 1 hour. The AAL3 paragraph (chunk 31) was at rank 4, just outside the top 3, so Answer 1 gave the AAL2 value and misstated the AAL3 rule. The Cross Encoder moved chunk 31 to rank 1, with a score of 6.79 against 6.43 for the best AAL2 chunk, and Answer 2 said 15 minutes. Query 6 shows the same problem with a smaller margin. The top six baseline candidates (cosine 0.731 to 0.771) all came from the sections on assurance levels and reauthentication, and the Cross Encoder put the AAL1 chunk first with a score of 7.18 against 7.12 for the AAL3 chunk.

\textbf{Other retrieval fixes.} In query 1 the baseline top 3 were about creating and changing passwords (a blocklist check and two usability sections). The rule sat at rank 5, and Answer 1 only pointed to Sec. 3.1.1 without stating it. In query 3 two of the three baseline chunks came from Appendix B on syncable authenticators and the third from the section on exportable keys at AAL2 or below. The sentence that forbids synced passkeys at AAL3 was at rank 7. In query 8 the top baseline chunk was a usability section that mentions SMS only as an example of a suitability problem, and the other two were also usability chunks. The chunk that calls PSTN the one restricted authenticator was at rank 4.

\textbf{Better retrieval was not always enough.} In query 2 the gold chunk moved from rank 10 to rank 1, but the LLM still said encryption is required and dropped the ``unless sent via the PSTN'' exception, although the complete sentence was in Passage 1. In query 4 the rule was at rank 1 and its first sentence was at the start of Passage 1, but the LLM answered ``Not found''. A likely reason is that the question says ``every 90 days'' while the rule says ``periodically'', but this was not tested. In both cases retrieval was fixed and the 1.5B LLM was the weak step.

\textbf{Counter examples.} In query 9 the Cross Encoder moved the rule ``Biometrics SHALL only be used as part of multifactor authentication'' (chunk 121) from rank 3 to rank 7. Its top 3 were a usability table, the introduction of the biometrics section and a chunk about problems collecting fingerprints. Its best score was minus 7.62 and all ten scores were at or below that, so it was not confident about any candidate. In query 11 the gold chunk (chunk 69, the same chunk that answers query 2) was not in the top 10, so reranking could not help. The Cross Encoder ranked first a chunk about issued recovery codes sent to a recovery address (postal, email, text message or voice), and Answer 2 said yes, so a safe ``Not found'' became a wrong answer.

\textbf{Wider test.} Retrieval alone was also run on 48 questions. These are all the questions that were tried while looking for failure cases, so they are not a random sample and the scores should not be read as typical performance.

\begin{center}
\small
\begin{tabular}{@{}lccc@{}}
\toprule
 & Hit@1 & Hit@3 & MRR@10 \\
\midrule
Pipeline 1, Bi Encoder & 0.250 & 0.583 & 0.437 \\
Pipeline 2, plus Cross Encoder & 0.500 & 0.750 & 0.617 \\
\bottomrule
\end{tabular}
\end{center}

The gold rank improved in 19 questions, stayed the same in 15 and became worse in 5. In the remaining 9 questions (19 percent) the gold chunk was never in the top 10, so Hit@10 is 0.812 for both pipelines. Five of these nine gold sentences are SHALL NOT rules (composition rules, password hints, voice biometrics, email for out of band codes, and manual entry codes not being phishing resistant). Reranking also helped with 100 word chunks (Hit@3 0.542 to 0.688) and 150 word chunks (0.583 to 0.708), so the result does not depend on the chunk size.

\section{Latency profiling}

Each step was timed with \texttt{time.perf\_counter()}, using 20 runs for each of the 11 queries after a warm up, on a laptop CPU (Ryzen 5 5500U, no GPU). The LLM time was measured once per answer.

\begin{center}
\small
\begin{tabular}{@{}lrrr@{}}
\toprule
Step & Mean & Median & 95th percentile \\
\midrule
Pipeline 1: Bi Encoder retrieval (top 10) & 17.8 ms & 17.5 ms & 23.2 ms \\
Pipeline 2: extra Cross Encoder step (10 pairs) & 365.7 ms & 350.4 ms & 489.2 ms \\
Pipeline 2: total retrieval & 383.5 ms & 368.1 ms & 509.4 ms \\
LLM generation, Answer 1 & 51.0 s & 50.5 s & 60.2 s \\
LLM generation, Answer 2 & 51.7 s & 51.9 s & 56.4 s \\
End to end, Pipeline 1 & 51.02 s & 50.51 s & 60.19 s \\
End to end, Pipeline 2 & 52.08 s & 52.25 s & 56.82 s \\
\bottomrule
\end{tabular}
\end{center}

Reranking raises retrieval time from 17.8 ms to 383.5 ms, about 21.5 times slower, because the Cross Encoder runs the transformer on 10 pairs while the Bi Encoder only encodes the question. The extra 366 ms is 0.7 percent of the baseline LLM time of about 51 seconds. The measured end to end time rose from 51.02 s to 52.08 s, about 2.1 percent. Reranking explains 0.37 s of this 1.06 s rise, and the other 0.70 s is the difference between the two LLM rows. LLM times varied a lot between questions (for Answer 1 the 95th percentile is 60.2 s against a median of 50.5 s) and each answer was timed once, so a 0.7 s difference is inside this spread and the data cannot assign it to reranking.

\section{Scaling question}

The Cross Encoder is not run over the whole collection because it needs the question and the chunk together, so nothing can be computed before the question arrives. For N chunks it runs the transformer N times for every question, and the cost grows linearly with N. This was measured for query 1: scoring all 413 chunks took 16.9 seconds (about 41 ms per pair, slightly higher than the 37 ms per pair above, from a separate run), which is 44 times slower than the two stage pipeline at 0.38 seconds. For this query the gold chunk was ranked first in both cases, so the extra time gave no better result, although one query is not enough to say this holds in general.

\begin{center}
\small
\begin{tabular}{@{}rrr@{}}
\toprule
Number of chunks & Cross Encoder on every chunk & Two stage pipeline \\
\midrule
413 (measured) & 16.9 s & 0.38 s \\
10,000 (estimated) & 6.8 minutes & 0.38 s \\
1,000,000 (estimated) & 11.3 hours & 0.38 s \\
\bottomrule
\end{tabular}
\end{center}

The Bi Encoder computes every chunk vector once, offline. At question time it encodes only the question and searches the stored vectors, and an approximate index such as HNSW keeps this in milliseconds even for millions of chunks. The Cross Encoder then sees only 10 candidates, so its cost stays the same however large the collection grows.

\section{Conclusion and limitations}

Reranking improved which chunks reached the LLM in every failure query and raised Hit@3 on the 48 questions from 0.583 to 0.750, for about 0.37 seconds per question. The gains appeared where the baseline faced several similar candidates (the AAL levels) or where the answer was a negation or an exception. It has clear limits. The gold rank became worse in 5 of 48 questions, it cannot recover the 9 questions where stage 1 missed the gold chunk, and the LLM still failed in queries 2 and 4 with the right text in front of it. Next I would retrieve more candidates (top 30 to 50), add keyword search (BM25) for exact terms such as ``email'' and ``PSTN'', and use a stronger LLM.

The limitations are that the 48 questions and gold sentences were written for this assignment and not taken from a benchmark, answers were judged by hand, the LLM is small, each LLM time was measured once, the full collection scan was timed on one query, and all timings come from one laptop.

\end{document}
